\item[\textbf{Exercise 15.3-3}]
What is an optimal Huffman code for the following set of frequencies, based on
the first 8 Fibonacci numbers?

\texttt{a:1 b:1 c:2 d:3 e:5 f:8 g:13 h:21}

Can you generalize your answer to find the optimal code when the frequencies are
the first $n$ Fibonacci numbers?

\Solution

\HuffmanStep
    {a:1 b:1 c:2 d:3 e:5 f:8 g:13 h:21}
    {exercise-15.3-3-step-1.pdf}
    {Step 1}

\HuffmanStep
    {c:2 a-b:2 d:3 e:5 f:8 g:13 h:21}
    {exercise-15.3-3-step-2.pdf}
    {Step 2}

\HuffmanStep
    {d:3 a-b-c:4 e:5 f:8 g:13 h:21}
    {exercise-15.3-3-step-3.pdf}
    {Step 3}

\HuffmanStep
    {a-b-c-d:7 e:5 f:8 g:13 h:21}
    {exercise-15.3-3-step-4.pdf}
    {Step 4}

\HuffmanStep
    {a-b-c-d-e:12 f:8 g:13 h:21}
    {exercise-15.3-3-step-5.pdf}
    {Step 5}

\HuffmanStep
    {a-b-c-d-e-f:20 g:13 h:21}
    {exercise-15.3-3-step-6.pdf}
    {Step 6}

\HuffmanStep
    {a-b-c-d-e-f-g:33 h:21}
    {exercise-15.3-3-step-7.pdf}
    {Step 7}


\subsection*{Huffman Code Table}

\begin{tabularx}{\linewidth}{ C | C | C }
Message & Frequency & Code \\
\hline
a & 1 & 1111110 \\
\hline
b & 1 & 1111111 \\
\hline
c & 2 & 111110 \\
\hline
d & 3 & 11110 \\
\hline
e & 5 & 1110 \\
\hline
f & 8 & 110 \\
\hline
g & 13 & 10 \\
\hline
h & 21 & 0 \\
\end{tabularx}

\subsection*{General Solution}

\noindent To find the optimal code when the frequencies are the first $n$
Fibonacci numbers, we can use the following algorithm:
\begin{enumerate}
    \item Generate the first $n$ Fibonacci numbers.
    
    \item Sort the frequencies in decreasing order.
    
    \item Initialize a priority queue with the characters and their
    corresponding frequencies.
    
    \item For $i = 1$ to $n - 1$:
    \begin{enumerate}
        \item Remove the two characters with the lowest frequencies from the
        priority queue.
        
        \item Create a new node with these two characters as children and with a
        frequency equal to the sum of their frequencies.
        
        \item Insert the new node back into the priority queue.
    \end{enumerate}
    
    \item Traverse the Huffman tree to assign codes to each character.
    
    \item Return the code table.
\end{enumerate}
